% 题干转写：完成
% 解析转写：完成
% 图形核对：完成（原题无题图）
\section{第37届复赛第4题：回旋加速器与反质子产生阈能}

\subsection{题干}
劳伦斯（E. O. Lawrence）在1930年首次提出回旋加速器原理：带电粒子在两个半圆形磁场中沿圆弧轨道旋转，反复通过两半圆缝隙间的高频电场而获得较高能量。

（1）费米实验室的高能质子同步加速器 Tevatron 的粒子运行最大回旋圆轨道周长为 $L_{\max}=6436\,\mathrm m$，可将质子加速到最大总能量 $E_{p,\max}=1.00\times10^6\,\mathrm{MeV}$。假设质子始终在垂直于匀强磁场的平面内运动，忽略电磁辐射引起的能量损失。求所需磁感应强度的最小值 $B_{\min}$。

（2）高能入射质子轰击静止靶质子可以产生反质子，反应为
\[
p+p\longrightarrow p+p+p+\bar p.
\]
求产生反质子时入射质子的最小动能，并判断 Tevatron 加速的质子能否产生反质子。已知质子质量 $m_p=938.3\,\mathrm{MeV}/c^2$，真空光速 $c=3.00\times10^8\,\mathrm{m/s}$。

\subsection{参考解答整理}
\paragraph{（1）磁场强度下限。}
相对论质子在垂直匀强磁场中满足
\[
\frac{\mathrm d\boldsymbol p}{\mathrm dt}=e\boldsymbol v\times\boldsymbol B,
\qquad
B=\frac{p}{eR}.
\]
最大轨道周长给出
\[
R\leqslant\frac{L_{\max}}{2\pi}.
\]
在给定动量下，半径取最大值时所需磁场最小。由能量--动量关系
\[
E_p^2=p^2c^2+m_p^2c^4,
\]
可得
\[
\boxed{B_{\min}=\frac{2\pi}{eL_{\max}c}
\sqrt{E_{p,\max}^2-m_p^2c^4}=3.26\,\mathrm T.}
\]

\paragraph{（2）固定靶反应阈能。}
设实验室系中入射质子的总能量、动量为 $E_1$、$p_1$。反应前系统的洛伦兹不变质量为
\[
M^2=\left(\frac{E_1+m_pc^2}{c^2}\right)^2-\left(\frac{p_1}{c}\right)^2
=\frac{2m_pE_1}{c^2}+2m_p^2.
\]
阈值时末态四个粒子在质心系中没有相对动能，因此
\[
M^2=(4m_p)^2.
\]
联立得到入射质子的最小总能量
\[
E_{1,\min}=7m_pc^2=6568.1\,\mathrm{MeV}.
\]
故其最小动能为
\[
\boxed{K_{1,\min}=E_{1,\min}-m_pc^2
=6m_pc^2=5629.8\,\mathrm{MeV}.}
\]
由于
\[
1.00\times10^6\,\mathrm{MeV}>6568.1\,\mathrm{MeV},
\]
Tevatron 加速的质子能够轰击静止靶质子产生反质子。

\subsection{图件与转写依据}
原题无独立题图，故不添加非原题的TikZ图件。题干、数值、公式与阈能均依据本题的题面转写稿和答案公式核对表；本题第（1）问22分，第（2）问18分。
